\documentclass[../main.tex]{subfiles}
\begin{document}

\section{Overview}

The previous chapter presented the proposed a novel summarization pipeline. This chapter evaluates whether the proposed summarization pipeline preserves the
causal structure of attribution graphs while making them smaller and easier to
interpret. The evaluation is organized around three questions. First, does
influence--relevance pruning retain enough of the replacement-model computation
to remain faithful to the original attribution graph? Second, does the proposed
ILP clusterer produce supernodes whose members have coherent downstream causal
effects? Third, do the final summary nodes and edges predict the behavior of the
model under intervention?

The chapter is therefore organized by experiment rather than by separating all
metrics, baselines, and results into different sections. Section~\ref{sec:setup}
describes the common datasets, implementation environment, and default pipeline
configuration. Section~\ref{sec:pruning-experiment} evaluates pruning.
Section~\ref{sec:clustering-experiment} evaluates supernode clustering.
Section~\ref{sec:steering-experiment} validates summary graphs with steering
interventions. Section~\ref{sec:multihop-stress-test} explains how the multi-hop
graphs are used as a qualitative stress test.

The main quantitative results support the three claims. Adding upstream
relevance to the pruning score preserves graph completeness and replacement
score at matched graph sizes. The ILP clusterer achieves the strongest
structural preservation among the clustering baselines: at the $p65$ role-cosine
threshold it obtains an internalized feature-edge mass of $0.0994$ and a DAG
loss of $0.0359$, while retaining a high role gap. Finally, the final summary
graphs obtain a sign consistency of $0.9361$ and a
Spearman edge-effect correlation of $0.5633$, indicating that their signed edges
largely agree with measured intervention responses.

\section{Common Experimental Setup}
\label{sec:setup}
\label{sec:datasets}

\subsection{Datasets}

\paragraph{Synthetic analogies from BATS.}
The main quantitative evaluation uses synthetic analogy prompts built from the
Bigger Analogy Test Set (BATS). BATS is a benchmark for probing relational and
semantic knowledge in word embeddings and language models. We use entities from
its lexicography category and construct prompts of the form ``The saying goes:
$A$ is to $B$ as $C$ is to $D$'', where $A$, $B$, $C$, and $D$ come from the
same subcategory. The model must predict $D$ from the relation expressed by
$A \to B$ and the query entity $C$.

The dataset contains 100 prompts in total, with 10 prompts for each of the
following relation types: country--capital, country--language, UK city--county,
name--nationality, name--occupation, animal--young, animal--sound,
animal--shelter, object--color, and male--female. The analogy setting is useful
because the answer entity $D$ should be recoverable by applying the relation
between $A$ and $B$ to $C$. A good summary should therefore preserve features
that carry the source entities, the relation, the query entity, and the predicted
answer. For the quantitative pruning and clustering experiments, we keep prompts
that the model answers correctly with confidence above $20\%$. This filter
ensures that every attribution graph explains a prediction the model actually
makes.

\paragraph{Multi-hop reasoning.}
The second dataset follows the multi-hop queries of Ameisen et
al.~\cite{ameisen2025circuit}. Each prompt requires several reasoning steps
before reaching the answer, so it tests whether the summary pipeline can preserve
causal structure beyond a single local association. The dataset contains 35
prompts covering basic knowledge in geography, biology, physics, and common
knowledge. These graphs are used as a qualitative stress test for the
summarization pipeline and viewer, while the quantitative tables focus on the
BATS-derived analogy graphs, where enough examples are available for controlled
sweeps.

\subsection{Environment and implementation}

Experiments run on a single workstation with an NVIDIA GeForce RTX 4090 GPU
(24\,GB VRAM) and 128\,GB of RAM, under Ubuntu 22.04 LTS. The implementation
uses Python 3.13.5. Attribution graphs are generated with the circuit-tracing
library~\cite{circuit-tracer}. The summary graphs reuse the same feature
metadata and are rendered with an adapted version of the library's user
interface

\subsection{Default pipeline configuration}

Unless stated otherwise, pruning uses the influence--relevance importance score
from Chapter~3 with an edge retained-mass threshold of $0.95$. The pruning sweep
varies the node retained-mass threshold and the mixture parameter $\alpha$ in
Eq.~\eqref{eq:importance}. The clustering experiments use pruned analogy graphs
from the $\alpha=0.5$ condition, which balances downstream influence and
upstream relevance. The exact clusterer is the binary integer linear program
(ILP) of Eq.~\eqref{eq:Lstage2}, solved through SciPy's \texttt{milp} interface
to the HiGHS mixed-integer solver with an explicit time limit. Thus the
mathematical model is an ILP, while the software backend is a mixed-integer
solver.

\section{Experiment 1: Pruning Faithfulness}
\label{sec:pruning-experiment}
\label{sec:eval-method}

\subsection{Question and comparison}

The pruning experiment asks whether the proposed influence--relevance score can
remove graph mass while preserving the end-to-end computation represented by the
attribution graph. The direct baseline is influence-only pruning from Ameisen et
al.~\cite{ameisen2025circuit}. It sets $\alpha=1$ in Eq.~\eqref{eq:importance},
so a node is ranked only by downstream reach to the target logit. The proposed
setting uses smaller values of $\alpha$ to combine downstream influence with
upstream relevance to input tokens. This comparison tests whether relevance
improves the semantic focus of the retained graph without sacrificing graph-level
faithfulness.

\subsection{Metrics}

Pruning is evaluated with two graph-faithfulness metrics. The first metric tests
whether influence can still be traced from the output back to retained input
tokens. The second metric tests how completely retained nodes explain incoming
edge mass throughout the graph.

\paragraph{Replacement score.}
Replacement score measures the fraction of end-to-end influence from input
tokens to output logits that flows through the retained part of the graph. Let
$V'$ and $E'$ be the retained node and edge sets from the pruning stage. The
complement $V \setminus V'$ denotes attribution-graph nodes outside the pruned
graph. Let $W'$ be the adjacency after this
merge, and let $\mathrm{Inf}'_v$ be the influence score computed from $W'$ using
Eq.~\eqref{eq:inf}. The score is
\begin{equation}
\mathrm{Rep}
=
\frac{\sum_{v \in V' \cap V_{\text{emb}}} \mathrm{Inf}'_v}
     {\sum_{v \in V} \mathrm{Inf}'_v}.
\label{eq:metric-replacement}
\end{equation}
This is a strict score: it is high only when output influence can be traced all
the way back to input tokens through retained pathways, rather than being
absorbed by nodes outside the pruned graph.

\paragraph{Graph completeness.}
Graph completeness measures the fraction of incoming edges to all nodes that
originate from the pruned graph, weighted by each node's influence on the output.
As in Chapter~3, $(W')^{\text{in}}$ is column-normalized with zero columns left
at zero when a node has no incoming edge mass. Let
\[
q'_v =
\begin{cases}
\sum_{u \in V'} (W')^{\text{in}}_{uv}
= 1 - \sum_{u \in V \setminus V'} (W')^{\text{in}}_{uv},
& \text{if node } v \text{ has positive incoming edge mass in } W',\\
0, & \text{otherwise.}
\end{cases}
\]
This is the fraction of node $v$'s normalized incoming edge mass that comes from
retained nodes. The score is
\begin{equation}
\mathrm{Comp}
=
\frac{\sum_{v \in V} q'_v(\mathrm{Inf}'_v)}
     {\sum_{v \in V}(\mathrm{Inf}'_v)}.
\label{eq:metric-completeness}
\end{equation}
Unlike replacement score, completeness gives partial credit to nodes that are
mostly explained by retained features even when some influence remains assigned
to nodes outside the pruned graph.

\paragraph{Irrelevant feature count}
The irrelevant feature count measure the number of features that are not relevant to the information we care about.
Ameisen et al.~\cite{ameisen2025circuit} mention the "boring circuit" problem where a large number of active features and graph edges are responsible for basic, obvious role, like task
related features, punctuation, stopwords, format scalfolding, ... We usually care more about how the model derives the semantics required to output the next token rather than the syntactic features.
Therefore, this metric is used to measure the proportion of features that are not relevant to the information we care about.
For the analogies dataset, we label the token positions of the entities $A$, $B$, $C$, and the final ``to'' token as relevant, and the remaining positions as irrelevant.
After pruning, we report two counts of retained nodes at irrelevant positions: the number of cross-layer transcoder feature nodes and the number of input token (embedding) nodes. Lower values mean that pruning discards more of the syntactic scaffolding.
We prune all attribution graphs with a node retained-mass threshold of $0.02$ and compare settings.
\subsection{Results}

Table~\ref{tab:pruning-results} reports the pruning results averaged over the
100 analogy graphs. Influence-only pruning attains the highest graph
completeness ($0.604$) and replacement score ($0.437$), but it retains the most
irrelevant nodes: on average $2.29$ feature nodes and $9.00$ input-token nodes
lie on irrelevant positions. Adding upstream relevance ($\alpha<1$) removes
almost all of this scaffolding while losing at most $0.02$ completeness and
$0.02$ replacement score. The two normalizations trade off the feature and token
axes differently. Entmax, which assigns exactly zero weight to low-importance
tokens, drives the irrelevant token count to about $0.1$, effectively dropping
every irrelevant input node; softmax keeps a denser token weighting and
therefore retains $2.8$--$3.5$ irrelevant token nodes. On the feature axis both
normalizations are comparable, with fewer than $0.13$ irrelevant feature nodes.
These results support the claim that influence--relevance pruning preserves
graph faithfulness at matched thresholds while discarding features and tokens
that are not semantically relevant to the analogy. For the rest of the experiments, we use our settings with Entmax normalization and $\alpha=0.5$.

\begin{table}[ht]
\centering
\caption{Pruning results on 100 BATS analogy graphs at edge retained-mass
threshold $0.95$ and node retained-mass threshold $0.02$. Values are means
across graphs. Irrelevant feature and token counts are the numbers of retained
feature and input-token nodes at irrelevant positions (lower is better);
graph completeness and replacement score are the graph-faithfulness metrics of
Eq.~\eqref{eq:metric-completeness} and Eq.~\eqref{eq:metric-replacement} (higher
is better). The best value in each column is in bold and the second-best value
is underlined; the $\alpha=1.0$ baseline ignores token weights, so its
normalization is not applicable.}
\label{tab:pruning-results}
\resizebox{\linewidth}{!}{%
\begin{tabular}{llrrrrr}
\toprule
Method & Normalization & $\alpha$ & Irrelevant feature count $\downarrow$ & Irrelevant token count $\downarrow$ & Graph completeness $\uparrow$ & Replacement score $\uparrow$ \\
\midrule
Influence only & -- & 1.0 & 2.29 & 9.00 & \textbf{0.604} & \textbf{0.437} \\
Ours & Entmax 1.25 & 0.75 & 0.13 & \underline{0.12} & \underline{0.598} & \underline{0.425} \\
Ours & Entmax 1.25 & 0.5 & 0.12 & \textbf{0.11} & 0.589 & 0.421 \\
Ours & Softmax & 0.75 & \textbf{0.04} & 3.48 & 0.591 & 0.423 \\
Ours & Softmax & 0.5 & \underline{0.07} & 2.77 & 0.583 & 0.418 \\
\bottomrule
\end{tabular}
}
\end{table}

\section{Experiment 2: Causal Supernode Clustering}
\label{sec:clustering-experiment}
\label{sec:baselines}

\subsection{Question and setup}

The clustering experiment asks whether the proposed ILP clusterer forms
supernodes with coherent causal roles. The experiment uses 100 pruned analogy
graphs from the entmax $\alpha=0.5$ condition. For the ILP solver, we set the
role-cosine threshold to the $65$th percentile of allowed feature pairs, denoted
$p65$, and set the upper bound for the number of supernodes to $K_{\max}=20$.
All clustering methods are
matched to the same number of feature supernodes as the proposed summary on each
graph, so the comparison isolates partition quality rather than graph size. Let
$F = V'_{\text{feature}}$ be the retained feature nodes, let $f$ be the
partition from clustering, and let $r(i)$ be the role vector of
Eq.~\eqref{eq:role}. All edge-mass formulas use the retained edge set $E'$ and
the convention that $W_{uv}$ is the edge from source $u$ to target $v$.

\subsection{Baselines}

\paragraph{Spectral clustering on role vectors.}
Spectral clustering on role vectors uses the same role-vector representation as
the proposed clusterer, but replaces the ILP objective with a standard spectral
partitioning algorithm. The affinity is a rectified cosine similarity between
role vectors. This baseline tests whether a generic graph-free clustering method
can recover comparable supernodes from the same feature representation.

\paragraph{K-means on role vectors.}
K-means on role vectors first unit-normalizes each role vector and then minimizes
Euclidean within-cluster variance. After normalization, Euclidean distance is a
cosine-distance surrogate, so this is a strong geometric compactness baseline.
It tests whether tight role-vector clusters are sufficient for causal summary
quality.

\paragraph{Spectral clustering on adjacency.}
Spectral clustering on adjacency ignores the role vectors and instead uses the
symmetrized absolute pruned adjacency matrix as its affinity. It groups features
that are strongly connected in the retained graph. This baseline tests whether
connectivity alone is enough, without the incoming and outgoing role profiles
used by the proposed method.

\paragraph{Random same-size clusters.}
The random same-size baseline keeps the same feature-cluster size distribution
as the proposed summary but randomly assigns features to clusters. It controls
for the effect of cluster count and cluster sizes. Any method that does not beat
this baseline is not extracting causal or geometric structure from the graph.

\subsection{Metrics}

All clustering metrics are computed independently for each graph and then
averaged as unweighted graph means. Same-cluster feature pairs are always used
exactly; cross-cluster pairs are sampled with a fixed seed only when their count
exceeds the evaluation cap.

\paragraph{Role gap.}
Role gap measures whether full signed causal-role vectors are more similar
inside a supernode than across different supernodes. Let
$r_i=[v_i^{\text{out}};v_i^{\text{in}}]$ be the signed role vector from
Eq.~\eqref{eq:role}, and let
$\mathcal{Q}_{\text{same}}=\{(i,j): i,j\in F,\, i<j,\, f(i)=f(j)\}$ and
$\mathcal{Q}_{\text{diff}}=\{(i,j): i,j\in F,\, i<j,\, f(i)\ne f(j)\}$ be
same-cluster and cross-cluster feature pairs. The cosine uses
$\cos(x,y)=x^\top y/(||x||\,||y||+\epsilon)$ with $\epsilon=10^{-12}$. The score
is
\begin{equation}
\mathrm{RoleGap}
=
\overline{\cos(r_i,r_j)}_{(i,j)\in\mathcal{Q}_{\text{same}}}
-
\overline{\cos(r_i,r_j)}_{(i,j)\in\mathcal{Q}_{\text{diff}}}.
\label{eq:metric-role-gap}
\end{equation}
Higher role gap means that members of the same supernode have more similar
combined incoming and outgoing signed roles than features assigned to different
supernodes.

\paragraph{Internalized edge mass.}
Internalized edge mass measures how much fine-grained causal weight disappears
inside feature supernodes. Unlike a graph-wide edge fraction, this metric is
restricted to retained feature--feature edges and excludes embedding-to-feature
and feature-to-logit edges:
\begin{equation}
\mathrm{Internalized}
=
\frac{\sum_{(u,v)\in E'_{FF}:\, f(u)=f(v)} |W_{uv}|}
     {\sum_{(u,v)\in E'_{FF}} |W_{uv}|}.
\label{eq:metric-internalized}
\end{equation}
Lower internalized edge mass means that more retained causal structure remains
visible as edges between supernodes.

\paragraph{DAG loss.}
DAG loss measures the fraction of aggregated superedge mass that points backward
under the summary graph's topological order. Let $A_{C\to C'}$ be the raw
block-sum edge from supernode $C$ to supernode $C'$ before antiparallel merging
or backward-edge removal, and let $\sigma(C)$ be the final supernode order. The
backward set is
$R_\sigma=\{(C,C'): A_{C\to C'}\ne 0,\, \sigma(C)>\sigma(C')\}$. The score is
\begin{equation}
\mathrm{DAGLoss}
=
\frac{\sum_{(C,C')\in R_\sigma}|A_{C\to C'}|}
     {\sum_{C\ne C'}|A_{C\to C'}|+\epsilon}.
\label{eq:metric-dagloss}
\end{equation}
Lower DAG loss means the partition is more compatible with the final acyclic
summary graph.

\subsection{Results}

Table~\ref{tab:clustering-baseline-entmax} reports method means on entmax-pruned
analogy and multihop graphs. K-means obtains the largest role gap on both
datasets, which shows that direct geometric clustering can form compact
role-vector groups. However, those groups hide more feature--feature edge mass
and require more backward-edge removal than the proposed ILP summary when
converted into an acyclic summary graph.

The gain columns compare each method with the random same-size baseline.
Positive values are better. Since the random-baseline role gap is near zero, the
role-gain column reports the absolute difference $m-m_{\mathrm{rand}}$ rather
than a percentage. For internalized mass and DAG loss, gain is the percentage
reduction from the random baseline,
$100(m_{\mathrm{rand}}-m)/m_{\mathrm{rand}}$.

\begin{table}[ht]
\centering
\caption{Clustering evaluation on entmax-pruned analogy and multihop graphs.
Higher role gap indicates more coherent signed causal roles inside supernodes. Lower
internalized edge mass and lower DAG loss indicate better preservation of visible
causal structure. Gains are measured against the random same-size baseline. \textbf{Best} and \underline{second-best} values are highlighted for each metric}
\label{tab:clustering-baseline-entmax}
\resizebox{\linewidth}{!}{%
\begin{tabular}{lrrrrrrrrrrrr}
\toprule
& \multicolumn{6}{c}{Analogies (100)} & \multicolumn{6}{c}{Multihop (36)} \\
\cmidrule(lr){2-7}\cmidrule(lr){8-13}
Method & Role gap $\uparrow$ & Int. mass $\downarrow$ & DAG loss $\downarrow$ & Role $\Delta$ $\uparrow$ & Int. gain $\uparrow$ & DAG gain $\uparrow$ & Role gap $\uparrow$ & Int. mass $\downarrow$ & DAG loss $\downarrow$ & Role $\Delta$ $\uparrow$ & Int. gain $\uparrow$ & DAG gain $\uparrow$ \\
\midrule
Random same-size clusters & -0.012 & \underline{0.082} & 0.124 & $+0.000$ & \underline{$+0.0\%$} & $+0.0\%$ & 0.004 & \underline{0.062} & 0.091 & $+0.000$ & \underline{$+0.0\%$} & $+0.0\%$ \\
Spectral on adjacency & 0.108 & 0.187 & 0.114 & $+0.121$ & $-128.2\%$ & $+8.1\%$ & 0.081 & 0.118 & 0.076 & $+0.077$ & $-89.7\%$ & $+16.0\%$ \\
Spectral on role vectors & 0.562 & 0.097 & \underline{0.078} & $+0.575$ & $-18.1\%$ & \underline{$+37.4\%$} & 0.479 & 0.087 & \underline{0.051} & $+0.474$ & $-40.0\%$ & \underline{$+43.6\%$} \\
K-means on role vectors & \textbf{0.713} & 0.084 & 0.092 & \textbf{$+0.726$} & $-2.3\%$ & $+25.9\%$ & \textbf{0.612} & 0.068 & 0.058 & \textbf{$+0.608$} & $-9.1\%$ & $+36.0\%$ \\
\midrule
Ours ILP & \underline{0.675} & \textbf{0.048} & \textbf{0.056} & \underline{$+0.688$} & \textbf{$+41.1\%$} & \textbf{$+54.9\%$} & \underline{0.591} & \textbf{0.048} & \textbf{0.038} & \underline{$+0.586$} & \textbf{$+22.7\%$} & \textbf{$+58.5\%$} \\
\bottomrule
\end{tabular}
}
\end{table}

The edge-mass metrics identify the intended advantage of the ILP formulation.
The ILP row has the lowest internalized feature-edge mass and DAG loss on both
datasets. Relative to random same-size clusters, ILP reduces internalized mass by
$41.1\%$ on analogies and $22.7\%$ on multihop, and reduces DAG loss by
$54.9\%$ on analogies and $58.5\%$ on multihop. K-means gives the largest
role-gap improvement, but it increases internalized mass over the random
baseline on both datasets. The ILP therefore sacrifices some role-vector
compactness to preserve more visible causal structure in the final summary
graph.

\section{Experiment 3: Intervention Validation via Entity-swap}
\label{sec:steering-experiment}

\subsection{Question and setup}

The intervention experiment asks whether the summary graph identifies
supernodes that are not only interpretable after inspection, but also causally
usable. In the analogy task, an Output-role supernode should contain features
that promote the answer entity $D$ for a prompt of the form ``$A$ is to $B$ as
$C$ is to''. If this interpretation is correct, replacing the source prompt's
Output-role features with the corresponding features from a donor prompt should
move the model toward the donor answer.

The experiment uses ordered source--donor pairs from the BATS analogy dataset.
The source and donor prompts must come from the same relation type, so the
intervention changes the answer entity while keeping the relation fixed. Pairs
are skipped when the source and donor have the same target token or when either
summary graph lacks usable Output-role CLT features. This filtering keeps the
test focused on answer transfer rather than on degenerate pairs or missing
summary nodes.

\subsection{Intervention protocol}

Entity-swap steering evaluates whether the Output supernodes discovered by the
summary graph are causally meaningful. For each analogy relation, we pair a source
prompt with a donor prompt from the same relation. We suppress the source
Output supernode and inject the donor Output supernode at the corresponding
position. If the summary graph has identified an answer-carrying supernode, this
intervention should move the model prediction away from the source answer and
toward the donor answer.

Figure~\ref{fig:summary-steering} shows the summary graph for one animal--young prompt
and the corresponding intervention view after replacing features that promote
``kid'' with features that promote ``calf''.

\begin{figure}[H]
\centering
\begin{subfigure}{0.95\linewidth}
\centering
\includegraphics[width=\linewidth]{Figure/our_summary.png}
\caption{Summary graph produced for the prompt ``The saying goes: fox is to cub as goat is to''.}
\end{subfigure}

\vspace{0.4cm}

\begin{subfigure}{0.95\linewidth}
\centering
\includegraphics[width=\linewidth]{Figure/our_steering.png}
\caption{Intervention view for the same prompt. The model output changes after
swapping Output features that promote ``kid'' with features that promote
``calf''.}
\end{subfigure}
\caption{The summary graph and interventions of the prompt ``The saying goes:
fox is to cub as goat is to''.}
\label{fig:summary-steering}
\end{figure}

\subsection{Metrics}

The primary metric is top-1 hit rate. A pair is counted as a top-1 hit only when
the steered model's most likely next token is the donor target token. This is a
strict success criterion because the intervention must overcome the clean source
answer and make the donor answer the model's final prediction.

The secondary metrics measure weaker forms of transfer. Top-5 hit rate counts
whether the donor target appears among the five highest-probability next tokens.
The donor probability change $\Delta p_{\text{donor}}$ is the steered donor
target probability minus its clean value, and the source probability change
$\Delta p_{\text{source}}$ is defined analogously for the source target. These
metrics are useful because a summary supernode can increase the donor answer
evidence and suppress the source answer without always making the donor answer
the top-1 prediction.

\subsection{Results}

Table~\ref{tab:entity-swap-steering} reports entity-swap transfer for two BATS
relations with the source suppression factor fixed at $-2$. Each row aggregates
50 attempted source--donor pairs for each donor factor, and hit rates are
computed over all attempted pairs. The results show that Output-role supernodes
carry transferable answer information: across the four settings, top-1 hit
rates range from $0.1600$ to $0.4400$, and top-5 hit rates range from $0.2800$
to $0.6800$.

\begin{table}[ht]
\centering
\caption{Entity-swap supernode steering on 50 BATS analogy pairs per relation
with the source factor fixed at $-2$. Donor Output features are injected with
the grouped donor factor. Hit rates are over all attempted pairs, and
$\Delta p$ values are steered probabilities minus clean probabilities. The best
value for each relation and metric is shown in bold.}
\label{tab:entity-swap-steering}
\resizebox{\linewidth}{!}{%
\begin{tabular}{lrrrrrrrr}
\toprule
\multirow{2}{*}{Relation}
& \multicolumn{4}{c}{Donor factor $4$}
& \multicolumn{4}{c}{Donor factor $8$} \\
\cmidrule(lr){2-5}\cmidrule(lr){6-9}
& Top-1 hit & Top-5 hit & $\Delta p_{\text{donor}} \uparrow$ & $\Delta p_{\text{source}} \downarrow$
& Top-1 hit & Top-5 hit & $\Delta p_{\text{donor}} \uparrow$ & $\Delta p_{\text{source}} \downarrow$ \\
\midrule
capital--country & 0.1600 & 0.2800 & \textbf{0.0944} & -0.5741 & \textbf{0.2000} & \textbf{0.4000} & 0.0817 & \textbf{-0.6746} \\
person--nationality & \textbf{0.4400} & \textbf{0.6800} & \textbf{0.2188} & -0.5328 & 0.3600 & 0.6000 & 0.1617 & \textbf{-0.5747} \\
\bottomrule
\end{tabular}
}
\end{table}

The transfer rate depends strongly on the relation. For capital--country pairs,
the best setting obtains a top-1 hit rate of $0.2000$ and a top-5 hit rate of
$0.4000$. For person--nationality pairs, the best setting obtains a top-1 hit
rate of $0.4400$ and a top-5 hit rate of $0.6800$. The donor probability change
shows the same pattern: the largest $\Delta p_{\text{donor}}$ for
capital--country is $0.0944$, whereas the largest value for
person--nationality is $0.2188$.

The factor sweep also shows that stronger donor injection is not uniformly
better. Increasing the donor factor from $4$ to $8$ improves the
capital--country top-1 and top-5 hit rates, but it reduces the
person--nationality hit rates and $\Delta p_{\text{donor}}$. This behavior
suggests that the entity-swap intervention is sensitive to activation scale, so
the results should be interpreted as evidence for transferable Output-role
information rather than as an optimized steering method. Overall, the
intervention validates the summary graph at the supernode level: the same nodes
that receive answer-role labels can be used to move the model toward a different
answer entity.

\section{Chapter Summary}

This chapter evaluated the summarization pipeline through three experiments. The
pruning experiment showed that adding upstream relevance to the importance score
does not substantially reduce replacement score or graph completeness at matched
graph sizes, while discarding features that are not semantically relevant to the
input prompt. The clustering experiment showed that the ILP clusterer yields the
lowest internalized edge mass and DAG loss among the tested baselines, while
K-means yields the strongest role-vector gaps. The end-to-end steering
experiment validated the final summary graph's ability to predict the behavior
of the model under intervention. Taken
together, these results indicate that the proposed pipeline preserves important
causal structure while compressing attribution graphs into more interpretable
summary graphs. Chapter~6 discusses the remaining limitations, and future work.
\end{document}
